\section{Calculus of Variations by Gelfand \& Fomin}

\subsection{Elements of the Theory}
A functional is defined as a function where the variable is in it of itself a function. It can also be seen as a class of functions.

While the calculus of functionals most developed and used branch is finding the extrema of variations, this is hardly its extent.

A great example is using the path length formula $L = \int_a^b \sqrt{1+\dot{y}^2}dx$ and finding which y achieves the minimum value for L, this is by definition the straightest line, or geodesic. It can also attempt to maximize. It is most often applied to any form of optimization theory or minimizing the action for Lagrandian dynamics.

In this book they will primarily use geometric language defining a space, the functional space, to be a space whose elements are functions. The nature and dimensions of such a space are defined by the functional defined. It should be as specific and defined with the problem. Continuity is the main constraint. To formalize the "closeness" needed for continuity the space will use normed linear space. Which is defined by the following axioms:

\begin{itemize}
  \item $x+y = y+x$
  \item $(x+y)+z=x+(y+z)$
  \item There exists a 0 element
  \item $\forall x \exists (-x): -x+x=0$
  \item Multiplication of 1
  \item Multiplatice communication
\end{itemize}
Normed is defined to be the magnitude of the length.

As I am flipping through this book I am seeing that the subject matter is quite trivial but presented in a overtly complex way. I may return to it but that is unlikely.
